% !PRINTED 527
\weberpage{527}
Now, by a well-known theorem from the theory of the
$\Gamma$-functions, if $k$ denotes an arbitrary quantity with positive
 imaginary part:
\[ \frac{1}{(-ik)^s} = \frac{(2\pi)^s}{\Gamma(s)}\int_0^\infty e^{2\pi ik\xi}\,\xi^{s-1}\,d\xi, \]
in which the power $(-ik)^s$ is uniquely defined by the convention that, when
$-ik = \varrho e^{i\Theta}$ is set, $\varrho$ is taken positive and the angle $\Theta$ between $-\frac{\pi}{2}$
and $+\frac{\pi}{2}$, one is to set $(-ik)^s = \varrho^s e^{i\Theta s}$.

If we here replace $ik$ by the conjugate imaginary
quantity $-ik'$, at the same time choosing a new letter of integration
 $\eta$, and multiply the two equations, there follows:
\weq{7}{\frac{1}{(kk')^s} = \frac{(2\pi)^{2s}}{\Gamma(s)\Gamma(s)}\int_0^\infty\int_0^\infty e^{2\pi i(k\xi-k'\eta)}(\xi\eta)^{s-1}\,d\xi\,d\eta.}
In this we substitute for $k, k'$ the two conjugate factors
$x+\omega_1y$, $x-\omega_2y$ of $\psi$, and introduce the integral expression (7)
into (6).

In this way we obtain:
\weq{8}{M_s = \frac{(2\pi)^{2s}}{A_s\Gamma(s)\Gamma(s)}\Sigma^{x,y}2\int_0^\infty\int_0^\infty e^{2\pi i[x(\xi-\eta)+y(\omega_1\xi+\omega_2\eta)]}(\xi\eta)^{s-1}\,d\xi\,d\eta.}
% SOURCE-ORIGINAL-SIC: Equation (8) prints A_s in the denominator although equation (6) and the derivation indicate A^s; the source subscript is preserved.
We now regard the product of the two integrals on the
right-hand side of this expression, which for brevity we denote by
\weq{9}{W = 2\int_0^\infty\int_0^\infty e^{2\pi i[x(\xi-\eta)+y(\omega_1\xi+\omega_2\eta)]}(\xi\eta)^{s-1}\,d\xi\,d\eta}
as a double integral which, if $\xi, \eta$ are interpreted as rectangular
 coordinates in the plane, extends over the positive
quadrant of the coordinate system.

% Fig. 2 source description: axes eta and xi; the bisector separates regions I and II. At fixed u the hatched integration strips extend from v=u toward v=infinity, not from the origin.
\weberfigureimage{weber_b3_i0558_fig2_source_crop.png}{0.30}

In order to transform the double integral, we divide the domain of integration
 by a straight line bisecting the angle into the two
parts I, II (see Fig. 2), and substitute in the first part
%CONT>>
